\section{Left Factoring}
\label{sec:ch09-left-factoring}

\subsection{Masalah Faktor Umum}

Dua alternatif yang dimulai dengan simbol yang sama menyulitkan
parser top-down. Contoh:
\[
\text{stmt} \to \text{if}\ E\ \text{then}\ S\
\text{else}\ S \mid \text{if}\ E\ \text{then}\ S.
\]
Parser tidak dapat memilih alternatif hanya dengan melihat token
pertama, karena keduanya dimulai dengan \code{if}. Transformasi pada
bagian ini mengikuti Mogensen dan Thain; Juarna berhenti pada
penghapusan rekursi kiri.

\subsection{Transformasi}

Left factoring mengeluarkan faktor umum:
\[
A \to \alpha\beta_1 \mid \alpha\beta_2
\]
menjadi
\[
A \to \alpha A', \qquad A' \to \beta_1 \mid \beta_2.
\]
Jika salah satu alternatif berakhir tepat setelah \(\alpha\), tambahkan
\(\varepsilon\) pada \(A'\).

\begin{contoh}[If-else]
Aturan di atas menjadi:
\[
\begin{aligned}
\text{stmt} &\to \text{if}\ E\ \text{then}\ S\ \text{stmt}'\\
\text{stmt}' &\to \text{else}\ S \mid \varepsilon.
\end{aligned}
\]
Sekarang parser dapat memutuskan setelah membaca bagian
\code{then S}: jika token berikutnya \code{else}, pilih alternatif
pertama; jika tidak, pilih \(\varepsilon\).
\end{contoh}

\begin{contoh}[Identifier dan pemanggilan]
Aturan
\[
\text{primary} \to \code{ID} \mid \code{ID}\ \text{"("}\
\text{args}\ \text{")"}
\]
difaktorkan menjadi
\[
\begin{aligned}
\text{primary} &\to \code{ID}\ \text{primary}'\\
\text{primary}' &\to \text{"("}\ \text{args}\ \text{")"}
\mid \varepsilon.
\end{aligned}
\]
Dengan bentuk ini, parser membaca \code{ID} sekali, lalu memutuskan
berdasarkan token berikutnya.
\end{contoh}

\begin{figure}[H]
\centering
\begin{tikzpicture}[node distance=6mm and 10mm]
  \node[phase, text width=4.6cm] (before)
        {Sebelum: \(\alpha\beta_1 \mid \alpha\beta_2\)};
  \node[phase, text width=4.6cm, right=of before] (after)
        {Sesudah: \(\alpha A'\)};
  \node[phase, text width=4.6cm, below=of after] (tail)
        {\(A' \to \beta_1 \mid \beta_2\)};
  \draw[arrow] (before) -- (after);
  \draw[arrow] (after) -- (tail);
\end{tikzpicture}
\caption{Left factoring mengeluarkan faktor umum}
\label{fig:ch09-left-factoring}
\end{figure}

\subsection{Kapan Left Factoring Diperlukan}

Left factoring diperlukan jika dua alternatif memiliki prefiks umum
yang tidak dapat dibedakan dengan satu token lookahead. Pada grammar
Subset-C, faktor umum muncul pada:
\begin{itemize}[leftmargin=*]
  \item \code{primary}: \code{ID} versus \code{ID ( args )};
  \item \code{stmt}: \code{ID = expr ;} versus \code{ID ( args ) ;};
  \item \code{topdecl}: deklarasi fungsi dengan dan tanpa tubuh.
\end{itemize}

Setelah factoring, grammar lebih cocok untuk parser LL(1). Namun,
factoring dapat membuat grammar lebih sulit dibaca manusia. Karena
itu, grammar asli sering dipertahankan untuk dokumentasi, sedangkan
grammar hasil transformasi dipakai untuk implementasi.

\begin{warningbox}
Left factoring tidak menghilangkan ambiguitas. Ia hanya mengurangi
konflik FIRST. Dangling else tetap memerlukan aturan resolusi
tambahan meskipun grammar sudah difaktorkan.
\end{warningbox}
